\section{Routh reduction by stages}\label{section6}
In reduction by stages, we study the reduction of a $G$-invariant system (symplectic or Lagrangian) under the action of the full group $G$ and under the induced action w.r.t.\ a normal subgroup $K\lhd G$. We shall adopt as far as possible the notations used in~\cite{MarsdenHamRed}. A detailed construction of the following def\/initions is found in this reference.
\begin{definition}\mbox{} \label{def:notations}
\begin{enumerate}\itemsep=0pt
\item The Lie-algebra of $K$ is $\lak$ and $i$ denotes the injection $i: \lak\to \lag$ with dual $i^*: \lag^*\to \lak^*$.
\item The group $G$ acts on $\lak$ by restriction of the adjoint action. The induced action of $G$ on $\lak^*$ is denoted by the same symbol ${\rm Ad}^*:G\times\lak^*\to \lak^*$.
\item  $\mu$ denotes an element in $\lag^*$ and $\nu\in\lak^*$. Then $G_\mu$ is the isotropy subgroup of $\mu$ under the ${\rm Ad}^*$-action of $G$ on $\lag^*$; $G_\nu$ is the isotropy subgroup of $\nu$ under the ${\rm Ad}^*$-action of $G$ on $\lak^*$ obtained as the dual of the restricted $Ad$-action of $G$ on $\lak$; and $K_\nu$ is the isotropy of $\nu$ w.r.t.\ to standard coadjoint action of $K$ on $\lak^*$. These groups satisfy $G_\nu\cap K= K_\nu$ and~$K_\nu$ is normal in $G_\nu$.
\item $\lag_\nu$ and $\lak_\nu$ denote the Lie algebras of $G_\nu$ and $K_\nu$ respectively. $\bar G_\nu$ denotes the quotient group $G_\nu/K_\nu$ and its Lie algebra equals $\bar\lag_\nu= \lag_\nu/\lak_\nu$.
\item The projections onto the quotient groups are denoted by $r: G\to \bar G=G/K$ and $r_\nu:G_\nu\to\bar G_\nu$, and on the level of the Lie algebra: $r': \lag \to \bar \lag =\lag/\lak$ and $r'_\nu:\lag_\nu \to \bar\lag_\nu$.   The inclusion map $G_\nu \to G$ induces a map $k_\nu: \lag_\nu \to \lag$, with its dual $k_\nu^*: \lag^*\to \lag_\nu^*$.
\item $\rho$ denotes an element in $\bar\lag_\nu^*$.
\end{enumerate}
\end{definition}

In~\cite{MarsdenHamRed} symplectic reduction by stages is performed under the condition of a so-called `stages hypothesis'. An element $\mu\in\lag^*$ is said to satisfy the stages hypothesis if for any $\mu'\in\lag^*$ satisfying $\mu|_\lak=\mu'|_\lak=\nu$ and $\mu|_{\lag_\nu} = \mu'|_{\lag_\nu} =\bar\nu$, there exists an element $k\in K_\nu$ and $g\in (G_\nu)_{\bar\nu}$ such that ${\rm Ad}^*_{kg} \mu'= \mu$.
The stages hypothesis is a condition on a chosen momentum value and depends only on the symmetry group $G$. It was already clear in~\cite{MarsdenHamRed} that the hypothesis is automatically satisf\/ied if $G$ is a central extension or if $G$ is a semi-direct product group. In the recent contribution~\cite{mikityuk} it has been pointed out that the hypothesis is in fact always satisf\/ied, and that it can be taken out of the reduction by stages statements altogether. Taking advantage of this result, in this paper, we will not make further reference to the stages hypothesis.


\subsection{Symplectic reduction by stages}\label{section6.1}

\begin{theorem}[Symplectic reduction by stages~\cite{MarsdenHamRed}]\label{thm:hamredstages}
Let $(M,\omega)$ be a symplectic manifold with a~canonical $G$-action $\Psi^M$ with an equivariant momentum map $J_G$.
\begin{enumerate}\itemsep=0pt
\item[$1.$] Fix a regular value $\mu\in\lag^*$ of the momentum map and perform symplectic reduction to obtain the symplectic manifold $(M_\mu,\omega_\mu)$.
\item[$2.$] The restriction of the action $\Psi^M$ to $K$ is canonical and the map $J_K=i^*\circ J_G:M\to \lak^*$ determines an equivariant momentum map for this induced action. Fix a regular value $\nu$ of $J_K$ and perform symplectic reduction to obtain the symplectic manifold $(M_\nu,\omega_\nu)$.


\item[$3.$] The level set $J^{-1}_K(\nu)$ is $G_\nu$-invariant.
\item[$4.$] The group $\bar G_\nu =G_\nu/K_\nu$ acts on $M_\nu$ by projecting the restricted action of $G_\nu$ on $J^{-1}_K(\nu)$. This induced action $\Psi^{M_\nu}$ is free, proper and canonical on $(M_\nu, \omega_\nu)$. Assume that $K_\nu$ is connected.
\item[$5.$] Fix an element $\bar \nu$ in $\lag_\nu^*$ such that the restriction of $\bar \nu|_{\lak_\nu}$ equals $\nu|_{\lak_\nu}$. There is a well-defined momentum map $J_{\bar G_\nu}: M_\nu \to \bar\lag^*_\nu$  for the induced action $\Psi^{M_\nu}$. This momentum map is determined from $J_G$ and $\bar \nu$ and has a non-equivariance cocycle:  $J_{\bar G_\nu}$ satisfies
\[(r'_\nu)^* \circ J_{\bar G_\nu} \circ \pi_\nu = k_\nu^*\circ J_G \circ i_\nu  -\bar \nu.\]
\item[$6.$] Let $\rho\in \bar\lag^*_\nu$ be a regular value for the momentum map $J_{\bar G_\nu}$ and let $(\bar G_{\nu})_\rho$ be the isotropy subgroup of $\rho$ w.r.t.\ the affine action of $\bar G_\nu$ on $\bar\lag_\nu^*$. Perform symplectic reduction to obtain the symplectic manifold $((M_\nu)_\rho, (\omega_\nu)_\rho)$ with $(M_\nu)_\rho = J^{-1}_{\bar G_\nu}(\rho)/(\bar G_\nu)_\rho$.
 \end{enumerate}
  If $\rho$ is chosen such that $(r'_\nu)^*\rho = \mu|_{\lag_\nu} - \bar \nu$, then there exists a symplectic diffeomorphism \[F: \ (M_\mu,\omega_\mu)\to ((M_\nu)_\rho, (\omega_\nu)_\rho).\]
\end{theorem}
For our purpose it is also important to understand the reduction of a $G$-invariant Hamiltonian~$h$ on~$M$. We assume that all conditions in Theorem~\ref{thm:hamredstages} are satisf\/ied. First note that, by def\/inition of the momentum maps, we have an inclusion $j_\mu$ of $J^{-1}_{G}(\mu)$ in $J^{-1}_K(\nu)$. Recall that we use $\pi_\nu$ for the projection $J^{-1}_K(\nu) \to M_\nu$. It was shown in~\cite{MarsdenHamRed} that the image of $\pi_\nu\circ j_\mu: J^{-1}_G(\mu) \to M_\nu$ is contained in $J^{-1}_{\bar G_\nu}(\rho)$. Moreover, this map is equivariant w.r.t.\ the action of $G_\mu$ on $J^{-1}_G(\mu)$ and $(\bar G_\nu)_\rho$ on $J^{-1}_{\bar G_\nu}(\rho)$ (this makes sense, since $G_\mu$ projects to a subset of~$(\bar G_\nu)_\rho$). The quotient of $\pi_\nu\circ j_\mu$ is the symplectic dif\/feomorphism $F$ mentioned in the previous theorem (see also Fig.~\ref{fig:diag5}).
\begin{figure}[htb]
\centering
\includegraphics{Langerock-Fig4}

\caption{Commuting diagram relating the dif\/ferent reduced symplectic manifolds.}\label{fig:diag5}\end{figure}

Let $h$ be a $G$-invariant Hamiltonian on $M$ and let $h_\mu$ be the function on $M_\mu$ obtained from $\pi_\mu^*h_\mu = i^*_\mu H$. On the other hand we let $h_\nu$ be the function satisfying $\pi^*_\nu h_\nu = i^*_\nu h$. This function is $\bar G_\nu$-invariant: $h_\nu([m]_{K_\nu} [g]_{K_\nu} ) = h(i_\nu(mg)) = h(i_\nu(m) g)$, with $m\in J^{-1}_K(\nu)$ and $g\in G_\nu$ arbitrary. Note that $(\pi_\nu\circ j_\mu)^* h_\nu = i^*_\mu h$.

The Hamiltonian $h_\nu$ is a $\bar G_\nu$-invariant function on $(M_\nu, \omega_\nu)$. Applying the second symplectic reduction to this manifold, we obtain a new reduced Hamiltonian $(h_\nu)_\rho$ on $(M_\nu)_\rho$.



\begin{proposition}\label{prop:hamred} $F^* ((h_\nu)_\rho) = h_\mu$.\end{proposition}
\begin{proof}
We rely on the commuting diagram in Fig.~\ref{fig:diag5}: \[\pi_\mu^*(F^* ((h_\nu)_\rho)) = (\pi_\nu\circ j_\mu)^* h_\nu = i^*_\mu (h).\] This uniquely characterizes $F^* ((h_\nu)_\rho)$ as the function $h_\mu$.
\end{proof}

\subsection{Routh reduction by stages}\label{section6.2}
Routh reduction by stages is symplectic reduction by stages applied to the symplectic structure of the initial Lagrangian system. In this section we show that the symplectic structures and energy hamiltonians in the dif\/ferent stages can in fact be associated to specif\/ic magnetic Lagrangians systems, and eventually gives us Routh reduction by stages. The symplectic reduction by stages then provides us a dif\/feomorphism relating the solutions of the dif\/ferent Euler--Lagrange equations for the Lagrangian systems in the f\/inal stages.

We start with a hyperregular Lagrangian $L$, invariant under the action of a Lie group $G$. We assume that this Lagrangian satisf\/ies a regularity condition which is more stringent than mere $G$-regularity.
\begin{definition} The Lagrangian $L$ is said to be $G$-hyperregular if for any $v_q\in TQ$ and any subspace $\lak' <\lag$ with injection $i':\lak'\to \lag$,  the mapping ${i'}^*\circ J_L|_{v_q}\circ i':\lak' \to \lak'^*$ def\/ined by $\xi \mapsto {i'}^*\big(J_L(v_q+(i'(\xi))_Q(q))\big)$ is invertible.
\end{definition}
Lagrangians of mechanical type are $G$-hyperregular. Let $K$ be a normal subgroup of $G$, $\lak$ the Lie algebra of $K$ and $i:\lak\to \lag$ the canonical injection. Due to the hyperregularity the invariant Lagrangian~$L$ is both $G$- and $K$-hyperregular, and both $G$- and $K$-invariant. By def\/inition of~$J_L$, the map $i^*\circ J_L= i^*\circ (\psi^{TQ})^*\circ \F L$ is the momentum map for the $K$-action.

\begin{theorem}[Routh reduction by stages]\label{thm:routhstages}
Assume $(Q,L)$ is a hyperregular, $G$-hyperregular and $G$-invariant Lagrangian system. Let $K$ denote a normal subgroup of $G$.
\begin{enumerate}\itemsep=0pt
\item[$1.$]%\label{stages:1}
Let $\mu\in\lag^*$ be  a regular value of the momentum map $J_L$ and $\A^0$ a $G$-connection on $Q$. Let $(Q/G_\mu \to Q/G,  L_0, \B_0)$ be the magnetic Lagrangian system obtained by performing Routh reduction with respect to $G$.
\item[$2.$]%\label{stages:2}
Fix a regular value $\nu\in \lak^*$ of the momentum map $i^*\circ J_L$ for the $K$-action and a $K$-connection $\A^1$ on $Q$. Assume that $\A^1$ is $G$-equivariant w.r.t.\ the action of $G$ on $Q$ and $\lak^*$. Consider the magnetic Lagrangian system  $(Q/K_\nu \to Q/K,  L_{1},  \B_{1})$ obtained by performing Routh reduction with respect to $K$.

\item[$3.$]%\label{stages:4}
 $\bar G_\nu$ acts on $Q/K_\nu$ and $Q/K$ by projecting the induced action of $G_\nu$ on $Q$. These induced actions are free and proper.
\item[$4.$]%\label{stages:5}
Assume that $K_\nu$ is connected. Fix an element $\bar \nu\in \lag^*_\nu$ such that $\bar \nu |_{\lak_\nu} = \nu|_{\lak_\nu}$. Then the magnetic Lagrangian system  $(Q/K_\nu \to Q/K,  L_{1},  \B_{1})$ is $\bar G_\nu$-invariant, $\bar G_\nu$-regular and admits a $\B_1\bar\lag_\nu$-potential $\delta_1$ entirely determined by the choice of $\bar \nu$. The potential satisfies, for arbitrary $q\in Q$
\[
(r'_\nu)^*\big(\delta_1([q]_{K_\nu})\big)=  -(\psi^Q\circ k_\nu)^*\big(%Ê
\langle \nu , \A^{1}(q)\rangle\big) + \overline\nu .
\] Let $J_1$ denote the momentum map associated with $\delta_1$.
\item[$5.$]%\label{stages:6}
Fix a regular value $\rho\in \bar\lag_\nu^*$ for the momentum map $J_{1}$ and let $(\bar G_{\nu})_\rho$ be the isotropy subgroup of $\rho$ w.r.t.\ the affine action of $\bar G_\nu$ on $\bar\lag_\nu$. Fix a $\bar G_\nu$-connection on $Q/K$. Consider the magnetic Lagrangian system $((Q/K_\nu)/(\bar G_\nu)_\rho \to (Q/K)/\bar G_\nu), L_{2}, \B_{2})$ obtained by performing Routh reduction with respected to $\bar{G}_\nu$.
%\item\label{stages:7} Assume that $\mu$ satisfies the stages hypothesis.
\end{enumerate}
If $\rho$ is chosen such that $(r'_\nu)^*\rho = \mu|_{\lag_\nu} - \bar \nu$, then every solution $\gamma(t)\in Q/G_\mu$ to the Euler--Lagrange equations for $(Q/G_\mu\to Q/G, L_{0} , \B_{0})$ is mapped to a solution in $(Q/K_\nu)/(\bar G_\nu)_\rho$ to the Euler--Lagrange equations for $((Q/K_\nu)/(\bar G_\nu)_\rho \to (Q/K)/\bar G_\nu, L_{2}, \B_{2})$. Conversely, a solution in $(Q/K_\nu)/(\bar G_\nu)_\rho$ to  the Euler--Lagrange equations for $((Q/K_\nu)/(\bar G_\nu)_\rho \to (Q/K)/\bar G_\nu, L_{2}, \B_{2})$ is the projection of a solution in $Q/G_\mu$ to the Euler--Lagrange equations for $(Q/G_\mu\to Q/G, L_{0} , \B_{0})$.
\end{theorem}
\begin{proof}
 1 and 2 are obtained by applying Routh reduction. 3~follows from~\cite[p.~152]{MarsdenHamRed}: we know that the quotient groups $\bar G_\nu=G_\nu/K_\nu$ acts in a~free and proper way on the quotient space $Q/K_\nu$.  The group $\bar G_\nu$ is a subgroup of $\bar G$ and acts freely and properly and $Q/K$. We now show~4.

{\bf $\bar G_\nu$-Invariance of the Routh reduced system $(Q/K_\nu\to Q/K, L_1,\B_1)$.}



\begin{lemma} If the connection $\A^{1}$ is chosen such that it is equivariant w.r.t.\ the action of the full group $G$, i.e.\ if
\[
(\Psi^Q_g)^*\A^{1} = {\rm Ad}_{g^{-1}}\A^{1},
\]then the magnetic Lagrangian system $(Q/K_\nu\to Q/K, L_1,\B_1)$ is $\bar G_\nu$-invariant and $\delta_1$ is a $\B_1\bar\lag_\nu$-potential.
\end{lemma}
\begin{proof}
  We f\/irst show that $L_1$ is $\bar G_\nu$-invariant. For that purpose, we choose an arbitrary $\bar g\in \bar G_\nu$ and let $g\in G_\nu$ be a representative. Similar we choose a point $(v_{[q]_K},[q]_{K_\nu})\in T_{Q/K_\nu}(Q/K)$ such that it is the projection of $v_q\in (i\circ J_L)^{-1}(\nu)\subset TQ$. By def\/inition of the quotient action on $T_{Q/K_\nu}(Q/K)$, the action of $\bar g$ on an element $(v_{[q]_K},[q]_{K_\nu})$ equals the projection of $v_qg$. We now check the invariance of $L_1$ at an arbitrary point in $T_{Q/K_\nu}(Q/K)$:
  \begin{gather*}
    L_1\left(\Psi^{T_{Q/K_\nu}(Q/K)}_{\bar g}\big(v_{[q]_K},[q]_{K_\nu}\big)\right)   = L(\Psi^{TQ}_g(v_q) )- \langle \nu,\A^{1}(qg)(T\Psi^{Q}_g(v_q) )\rangle \\ \phantom{L_1\left(\Psi^{T_{Q/K_\nu}(Q/K)}_{\bar g}\big(v_{[q]_K},[q]_{K_\nu}\big)\right)}{}
    = L(v_q) - \langle {\rm Ad}_{g^{-1}}^*\nu , \A^{1}(q)(v_q) \rangle  = L_1\left(v_{[q]_K},[q]_{K_\nu}\right).
  \end{gather*}
Next, we check the $\bar G_\nu$-invariance of $\B_1$.  Recall that $\B_1$ is the projection to $Q/K_\nu$ of the 2-form $d\langle\nu,\A^{1}\rangle$ on $Q$. We f\/irst consider the equivariance of this 2-form under $G_\nu$. Let $g\in G_\nu$ be arbitrary, then
  \[
  (\Psi^{Q}_{g})^*(d\langle\nu,\A^{1}\rangle) =  d\langle {\rm Ad}^*_{g^{-1}}\nu ,\A^1\rangle = d\langle\nu,\A^{1}\rangle.
  \]We thus obtain $G_\nu$-invariance for $d\langle\nu,\A^{1}\rangle$, and we may conclude that $(\Psi^{Q/K_\nu}_{\bar g})^* \B_1 = \B_1$ holds on $Q/K_\nu$.

  The third and f\/inal step is the def\/inition of the $\B_1\bar\lag_\nu$-potential. We consider an element $\bar \xi = [\xi]_{\lak_\nu} \in \bar\lag_\nu=\lag_\nu/\lak_\nu$ and let $\xi \in \lag_\nu$ be a representative. Then, by def\/inition of $\B_1$, the 1-form $i_{\bar \xi_{Q/K_\nu}}\B_1$ is the projection to $Q/K_\nu$ of the 1-form $i_{\xi_Q}d(\langle\nu,\A^{1}\rangle)$ on $Q$ (i.e.\ $\xi_Q$ projects to $\bar \xi_{Q/K_\nu}$).
  Again we concentrate on the 1-form on $Q$: \[
  i_{\xi_Q}d\langle\nu,\A^{1}\rangle = {\cal L}_{\xi_Q}(\langle\nu,\A^{1}\rangle) - d\big(i_{\xi_Q}\langle\nu,\A^{1}\rangle\big).\] Since $\langle\nu,\A^{1}\rangle$ is $G_\nu$-invariant, we conclude that $i_{\xi_Q}d\langle\nu,\A^{1}\rangle = - d\big(i_{\xi_Q}\langle\nu,\A^{1}\rangle\big)$. The exact 1-from on the right gives a strong hint of the structure of the $\bar\lag_\nu$-potential. Assume now that we f\/ixed an element $\overline\nu \in \lag_\nu^*$ such that $\overline\nu|_{\lak_\nu}=\nu|_{\lak_\nu}$.

  The function $\delta$  on $Q$,  def\/ined by
  \[-\delta_\xi(q) = \langle\nu,\A^{1}(q)(\xi_Q(q))\rangle - \langle \overline\nu,\xi\rangle \] is our candidate for the $\B_1\bar\lag_\nu$-potential. This statement makes sense provided that $\delta_\xi$ projects to a function on $Q/K_\nu$ and that it only depends on the equivalence class $\bar \xi=\xi+\lak_\nu$ of $\xi\in\lag_\nu$. The latter is a straightforward consequence of the fact that $\A^{1}$ is a principal $K$-connection. The $K_\nu$-invariance is more involved, and we rely on a result in~\cite{MarsdenHamRed}. For any $k$ in $K_\nu$, we have
  \begin{gather*}
  - \delta_{\xi}(qk)=  \langle \nu, \A^{1}(qk)(\xi_Q(qk))\rangle - \langle \overline\nu,\xi\rangle = \langle \nu, \A^{1}(q)(({\rm Ad}_{k}\xi)_Q(q))\rangle - \langle \overline\nu,\xi\rangle.
  \end{gather*}
  Therefore $\delta_\xi$ is constant on the orbits of $K_\nu$ in $Q$ if
  $\langle \nu, \A^{1}(q)((\xi-{\rm Ad}_{k}\xi)_Q(q))\rangle$ vanishes for all~$k$. To show this we introduce a function $f$ on $K_\nu$ given by $f(k)=\langle \nu, \A^{1}(q)((\xi-{\rm Ad}_{k}\xi)_Q(q))\rangle$ and we use similar arguments as in~\cite[p.~156]{MarsdenHamRed}. If we can show that $f(e)=0$,  $df|_e=0$ and $f(k_1k_2)=f(k_1)+f(k_2)$ for arbitrary $k_{1,2}\in K_\nu$, we may conclude that  $f=0$ (since $K_\nu$ is assumed connected).

The f\/irst condition $f(e)=0$ is trivial. To check the second condition: let $\kappa \in \lak_\nu$ be arbitrary, then
  \[
  df|_e(\kappa) = \langle \nu, \A^{1}(q)(-{\rm ad}_{\kappa}\xi)_Q(q))\rangle = -\langle \nu, {\rm ad}_\kappa\xi\rangle  = -\langle {\rm ad}^*_\kappa,\xi\rangle=0.
  \]
Above, we have used the fact that $K_\nu$ is normal in $G_\nu$ and that, as a consequence, the Lie bracket $[\kappa,\xi]$ is in $\lak_\nu$. Therefore the contraction of the corresponding fundamental vector f\/ield with $\A^{1}$ is precisely $[\kappa,\xi]$. Next, we check the third condition and compute $f(k_1k_2)$. Given the identity
  \[
  \xi-{\rm Ad}_{k_1}{\rm Ad}_{k_2} \xi =\xi -{\rm Ad}_{k_1}\xi+{\rm Ad}_{k_1}(\xi-{\rm Ad}_{k_2}\xi)
  \]
  and the fact that $k_1\in K_\nu$,
  \begin{gather*}
    f(k_1k_2) = \langle \nu, \A^{1}(q)((\xi-{\rm Ad}_{k_1k_2}\xi)_Q(q))\rangle\\
    \phantom{f(k_1k_2)}{} = \langle \nu, \A^{1}(q)((\xi-{\rm Ad}_{k_1}\xi)_Q(q)) \rangle + \langle \nu, \A^{1}(q)({\rm Ad}_{k_1}(\xi-{\rm Ad}_{k_2}\xi)_Q(q))\rangle\\
    \phantom{f(k_1k_2)}{} = f(k_1) +  \langle {\rm Ad}_{k_1}^*\nu, \A^{1}(q)((\xi-{\rm Ad}_{k_2}\xi)_Q(q))\rangle = f(k_1)+f(k_2).
  \end{gather*}
This completes the proof: the $\lag_\nu^*$-valued function $\delta$ is shown to be projectable to a $\bar\lag_\nu^*$-valued function on $Q/K_\nu$.  This is the sought-after potential $\delta_1$: for arbitrary $q\in Q$, we have \[
(r'_\nu)^*\big(\delta_1([q]_{K_\nu})\big)=  -(\psi^Q\circ k_\nu)^*\big(%Ê
\langle \nu , \A^{1}(q)\rangle\big) + \overline\nu . \tag*{\qed}
\]\renewcommand{\qed}{}
\end{proof}

{\bf Symplectic structure of $(Q/K_\nu\to Q/K, L_1,\B_1)$ and symplectic reduction by stages.}
\begin{lemma} Apply symplectic reduction by stages to the symplectic structure associated to the $G$-invariant Lagrangian system $(Q,L)$. Identify the symplectically reduced manifold $M_\nu$ with the symplectic structure on $T_{Q/K_\nu}(Q/K)$ induced by the magnetic Lagrangian system $(Q/K_\nu\to Q/K, L_1,\B_1)$. Then:
\begin{enumerate}\itemsep=0pt
\item[$1.$] The action $\Psi^{T_{Q/K_\nu}(Q/K)}$ of $\bar G_\nu$ on $T_{Q/K_\nu}(Q/K)$ is precisely the induced action on the first reduced space $M_\nu$ in symplectic reduction by stages.
\item[$2.$]  For a chosen $\bar \nu\in \lag_\nu^*$, the momentum map $J_{1}: T_{Q/K_{\nu}}(Q/K)\to \bar\lag_\nu^*$ associated with the magnetic Lagrangian system $(Q/K_\nu\to Q/K, L_1,\B_1)$ corresponds to the induced momentum map $J_{\bar G_\nu}$ from symplectic reduction by stages.
\end{enumerate}
\end{lemma}
\begin{proof}
1. The momentum map for the $K$ action is precisely $J_K:=i^*_\lak \circ J_L$, with $J_L:TQ\to \lag^*$. By def\/inition, the induced action of $\bar G_\nu$  on $J_K^{-1}(\nu)/K_\nu$ is obtained by projecting the action of~$G_\nu$ on~$J_K^{-1}(\nu)$. If we take into account that we realize the quotient manifold $J_K^{-1}(\nu)/K_\nu$ as $T_{Q/K_\nu}(Q/K)$, the induced action on $T_{Q/K_\nu}(Q/K)$ is obtained by projection of the action of $G_\nu$ on $TQ$ under the projection $TQ \to  T_{Q/K_\nu}(Q/K)$. This is precisely the action we have introduced above.

2. The {\em induced momentum map} ${J}_{\bar G_\nu}$ is def\/ined in the following way (we consider it directly as a function on $T_{Q/K_{\nu}}(Q/K)$ instead of on $M_\nu$):
\begin{gather}\label{eq:mommap}
\left\langle{J}_{\bar G_\nu}\big(v_{[q]_K}, [q]_{K_\nu}\big), \bar\xi\right\rangle = \langle J_L(v_q),k_\nu(\xi)\rangle - \langle \bar\nu,\xi\rangle = \langle \F L(v_q),\xi_Q\rangle - \langle \bar\nu,\xi\rangle,
\end{gather}
where $\xi \in \lag_\nu$ is arbitrary and projects to $\bar \xi \in \bar \lag_\nu$, $v_q$ projects to $(v_{[q]_K}, [q]_{K_\nu})$.  By def\/inition of the momentum map of the magnetic Lagrangian system $(Q/K_\nu\to Q/K, L_1,\B_1)$, we have
\begin{gather}\label{eq:mommap2}
\left\langle{J}_{1}\big(v_{[q]_K}, [q]_{K_\nu}\big), \bar\xi\right\rangle = \left\langle \F L_1\big(v_{[q]_K}, [q]_{K_\nu}\big),\big( (\bar \xi)_{Q/K},[p]_{K_\nu}\big)\right\rangle -  (\delta_1)_{\bar \xi}([p]_{K_\nu}).
\end{gather}
We now show that the right-hand side of~\eqref{eq:mommap2} equals the right-hand side of~\eqref{eq:mommap}. We therefore use the def\/inition of $\F L_1$ and $\delta_1$ as being the projection of maps upstairs:
\begin{gather*}
   \left\langle \F L_1\big(v_{[q]_K}, [q]_{K_\nu}\big),\big( (\bar \xi)_{Q/K},[q]_{K_\nu}\big)\right\rangle= \langle \F L(v_q),\xi_Q(q)\rangle - \langle \nu,\A^1(q)(\xi_Q(q))\rangle , \\
   (\delta_1)_{\bar \xi}([q]_{K_\nu})= \delta_\xi(p)= - \langle \nu,\A^1(q)(\xi_Q(q))\rangle + \langle \overline\nu,\xi\rangle. \tag*{\qed}
\end{gather*}\renewcommand{\qed}{}
\end{proof}

Finally, before we can reapply Routh reduction for the second stage, we need to check that $L_1$ is $\bar G_\nu$-regular.
\begin{lemma} The magnetic Lagrangian system $(Q/K_\nu\to Q/K, L_1,\B_1)$ is $\bar G_\nu$-regular.
\end{lemma}
\begin{proof}
  We have to show that, for any $(v_{[q]_K},[q]_{K_\nu})$ the map  \[\bar \lag_\nu \to \bar\lag_\nu^*; \qquad \bar \xi \mapsto {J}_{1}(v_{[q]_K}+\bar \xi_{Q/K}([q]_K),[q]_{K_\nu})\] is invertible. Let $v_q$ determine a tangent vector in $J^{-1}_K(\nu)$ representing $(v_{[q]_K},[q]_{K_\nu})$. Let $\eta$ denote an arbitrary element in $\bar \lag_\nu^*$. Due to the assumed $G$-hyperregularity, there is a unique $\xi\in\lag_\nu$ such that $k_\nu^*(J_L(v_q+ \xi_Q(q)))=(r'_\nu)^*\eta + \overline\nu$. The projection $\bar\xi=r'_\nu(\xi)$ of $\xi$ def\/ines the inverse element for $\eta$, since it is such that \[k^*_{\nu}(J_L(v_q+ \xi_Q(q))) - \bar\nu = (r'_\nu)^*{J}_{1} (v_{[q]_K}+\bar \xi_{Q/K})=(r'_\nu)^*\eta. \tag*{\qed}\] \renewcommand{\qed}{}
\end{proof}

{\bf Symplectic and Routh reduction by stages.}
Summarizing the above lemmas, we conclude that the magnetic Lagrangian system $(Q/K_\nu\to Q/K, L_1,\B_1)$ is amenable to Routh reduction and that the symplectic structure and momentum map associated to this Lagrangian system correspond to the symplectic structure and momentum map encountered in symplectic reduction by stages. If $\rho$ is chosen such that the compatibility relation $(r'_\nu)^*\rho = \mu|_{\lag_\nu} - \bar \nu$ holds, then from symplectic reduction by stages we have that the symplectic structures associated to $(Q/G_\mu\to Q/G, L_{0} , \B_{0})$ and $((Q/K_\nu)/(\bar G_\nu)_\rho \to (Q/K)/\bar G_\nu, L_{2}, \B_{2})$ are symplectically dif\/feomorphic by means of the symplectic dif\/feomorphism $F$ introduced earlier.  From  Proposition~\ref{prop:hamred} it follows that $F^*E_{L_2} = E_{L_0}$ and therefore the corresponding Hamiltonian vector f\/ields are $F$-related.  We def\/ine a map $\tau: Q/G_\mu \to (Q/K_\nu)/{(\bar G_\nu)_\rho}$ as $\tau([q]_{G_\mu}) = \big[[q]_{K_\nu}\big]_{(\bar G_\nu)_\rho}$. The map is well-def\/ined since $G_\mu$ is a subgroup of $G_\nu$ and since $r_\nu(G_\mu) \subset (\bar G_\nu)_\rho$.
\begin{figure}[htb]
\centering
\includegraphics{Langerock-Fig5}

\caption{Fibration of the symplectic dif\/feomorphism $F$.}\label{fig:diag4}
\end{figure}

\begin{lemma} The symplectic diffeomorphism $F$ is fibred over $\tau$.\end{lemma}
\begin{proof} We recall the def\/inition of the map $F$: f\/ix an element $(v_{[q]_G}, [q]_{G_\mu})$ and let $v_q\in J^{-1}_L(\mu)$ be a representative.  The point $F (v_{[q]_G}, [q]_{G_\mu})$ is obtained by taking the consecutive quotients of~$v_q$. In particular, the component of the f\/inal quotient in the conf\/iguration space $(Q/K_\nu)/(\bar G_\nu)$ of the magnetic Lagrangian system, is precisely the image $\tau$. \end{proof}

Since $F$ is a dif\/feomorphism, $\tau$ is onto. The Hamiltonian vector f\/ield on $T_{Q/G_\mu}(Q/G)$ and $T_{(Q/K_\nu)/(\bar G_\nu)_\rho}((Q/K)/\bar G_\nu)$ are $F$-related. Their integral curves project onto solutions of the Euler--Lagrange equations. This concludes the proof of Theorem~\ref{thm:routhstages}.
\end{proof}
